% ============================================================================
% Chapter 21 --- Wave 3: Network segmentation
% ----------------------------------------------------------------------------
% Purpose : declare zones, baseline ACL, Science DMZ established. Brownfield
%           trajectories dominate; greenfield is rare.
% Page target: ~6 pages.
% Dependencies: Wave 0 (Ch. 18); recommended Wave 2 (Ch. 20) to baseline
%               flows; reference layer Ch. 8 (zone catalogue D8.6, standard
%               interface discipline D8.7).
% Mirrors wave template (D18.2).
% ----------------------------------------------------------------------------
% Decisions registered in _coherence-EN.md:
%   D21.1 brownfield trajectory labels mapped to network reality
%   D21.2 cut-over per building/organizational-unit/segment, not institution-wide
%   D21.3 Science DMZ established in this wave
%   D21.4 exit criteria reference Network conformance subset (D16.5)
% ============================================================================

\chapter{Wave~3 --- Network segmentation}
\label{ch:wave3-network}

\begin{caveatbox}[Dependencies]
\noindent Wave~3 applies the wave chapter template established in
\trchap{ch:wave0-foundations}{}. Its pre-requisite is Wave~0;
Wave~2 is recommended (to baseline current flows before changing
the segmentation). It instantiates the network contract of
\trchap{ch:network}{}, in particular the zone catalogue of
\S\ref{sec:nw-reference} (D8.6) and the standard-interface
discipline of D8.7. Closure is gated by the four Network tests of
\trchap{ch:conformance}{}.
\end{caveatbox}

\paragraph{A note on the trajectory labels.} The trajectory labels
\emph{brownfield single-suite-centric} and \emph{brownfield hybrid}
are used consistently across the playbook, with their meaning
calibrated for identity, endpoint, and observability. In the
network layer, the operationally relevant distinction is between
single-vendor and multi-vendor estates rather than between
Microsoft and non-Microsoft. The labels are retained for
consistency; their meaning in this chapter is mapped accordingly.

%-----------------------------------------------------------------------------
\section{Goal and dependencies}
\label{sec:w3-goal}

The goal of Wave~3 is to bring the institution's network
infrastructure into conformance with the contract of
\trchap{ch:network}{}: a declared zone catalogue with documented
security profiles per zone, standards-aligned attachment via
RADIUS, declarative configuration via NETCONF/YANG (or equivalent),
flow telemetry in IPFIX, and a Science DMZ where the research
mission warrants it. The wave's success is measured against the
zone catalogue and the standard-interface discipline, not against
a complete refresh of the physical fabric.

\paragraph{What the wave does not undertake.} Wave~3 does not
require the institution to replace its network gear. It requires
that the gear be operated through standards on the operationally
relevant paths, that the zones be declared and enforced, and that
the conformance tests of \trchap{ch:conformance}{} pass on a
representative sample.

%-----------------------------------------------------------------------------
\section{Pre-requisites}
\label{sec:w3-prereq}

Beyond Wave~0, four operational conditions strengthen Wave~3:

\begin{itemize}[leftmargin=1.5em,nosep]
  \item a current inventory of network flows --- ideally
    established by Wave~2 --- so that the segmentation can be
    designed against actual usage rather than against assumed
    usage;
  \item an institutional addressing plan that supports the zone
    catalogue, with sufficient address space for the institution's planned
    growth (research, BYOD, IoT, or other drivers as applicable);
  \item a change-window agreement with the principal user
    constituencies (academic organizational units, research
    projects, administrative services), since per-segment cut-over
    requires scheduled
    downtime tolerable to those constituencies;
  \item an identified accountable team for the network operation,
    with authority to coordinate across organizational-unit IT
    functions where those exist independently.
\end{itemize}

%-----------------------------------------------------------------------------
\section{Coexistence pattern}
\label{sec:w3-coexistence}

Wave~3 is the wave with the most spatially distributed coexistence
in the playbook: the legacy flat network and the new zoned network
operate together for as long as the institution requires, with the
boundary of coexistence drawn at the building, the organizational
unit, or the network segment level rather than at the institutional
level.

\paragraph{Cut-over by perimeter.} Each building, organizational
unit, or segment is cut over to the new segmentation as a
self-contained
operation: the new zone is provisioned, the attachment service is
configured, the flow telemetry is activated, and the segment is
moved over. The remainder of the institution continues to operate
on the legacy flat network during this period. Per-segment cut-over
allows the institution to learn at modest scale before propagating
the design.

\paragraph{Bridge between zones during coexistence.} Some flows
must traverse the boundary during coexistence (services in a
migrated zone consumed by users in non-migrated buildings, for
example). The bridge is built on routing rules under change
control, with the explicit understanding that bridges are
temporary; bridges that have not been retired within an agreed
horizon are flagged for the governance body's attention.

\paragraph{Science DMZ established early.} The Science DMZ, where
present, is typically established early in the wave because its
construction is largely orthogonal to the rest of the segmentation.
Establishing it early demonstrates the architecture's research-
serving capability and decouples its delivery from the broader
cut-over schedule.

%-----------------------------------------------------------------------------
\section{Trajectory: greenfield}
\label{sec:w3-greenfield}

For institutions in a genuinely greenfield network position --- a
new building, a new site, a research-network refoundation --- the
trajectory is direct: the zone catalogue is implemented from the
start, the attachment service is the new RADIUS deployment, the
Science DMZ is a separate fabric, and the gear is procured against
the standard-interface contract. Indicative duration is six to
nine months for a new building of moderate size; longer for
site-scale greenfield, in which case the trajectory is split
into staged sub-projects.

%-----------------------------------------------------------------------------
\section{Trajectory: brownfield single-suite-centric}
\label{sec:w3-brownfield-ms}

For institutions whose network is built around a single dominant
vendor (a homogeneous Cisco, Juniper, Aruba, or Arista estate),
Wave~3's principal task is to operationalise the standard
interfaces of the existing gear without acquiring new equipment.

\paragraph{Pattern.} The existing gear is configured through
NETCONF/YANG (or its vendor-supported standard equivalent) for
the operationally relevant changes. The vendor's management
console is retained for diagnostic purposes only; production
changes flow through the standard interface under version control.
RADIUS attachment is reconfigured to consume identity attributes
from Wave~1's identity layer; flow telemetry is brought into the
collector layer of Wave~2.

\paragraph{What changes in Wave~3.} The zone catalogue is
declared, the attachment service is reconfigured for zone
assignment, the Science DMZ is established, the standard-interface
discipline is operationalised, and the conformance tests of
\S\ref{sec:cf-suite} pass on a representative sample. The gear
itself remains on its current refresh schedule; its eventual
substitution at refresh time is informed by the sovereignty grid
but not predetermined by Wave~3.

\paragraph{What does not change.} The physical topology, the
cabling, the gear inventory, and the existing vendor relationship.
The institutional procurement framework continues through this
wave; the wave produces information that the next procurement
cycle will use, not procurement decisions itself.

\paragraph{Indicative duration.} Twelve to twenty-four months,
dominated by the per-segment cut-over schedule and by the change
windows that user constituencies tolerate.

%-----------------------------------------------------------------------------
\section{Trajectory: brownfield hybrid (multi-vendor estate)}
\label{sec:w3-brownfield-hybrid}

For institutions whose network is multi-vendor --- often the
result of acquisitions, organizational-unit-level autonomy, or
research-network specialisation --- Wave~3's principal task is to
reconcile the
configuration model across vendors while preserving the existing
fabric.

\paragraph{Pattern.} The institution adopts a declarative
configuration management layer (for example an
infrastructure-automation toolchain or a source-of-truth-driven
configuration generator, or a comparable approach)
that emits NETCONF/YANG to each vendor and tolerates the residual
vendor-specific syntax for the cases where the standard does not
yet cover the configuration in question. The zone catalogue is
declared once at the institutional level and applied consistently
across vendors.

\paragraph{What changes in Wave~3.} The same zone catalogue and
attachment behaviour are realised across the multi-vendor estate;
flow telemetry is normalised at the collector layer (D12.3)
regardless of the vendor's IPFIX dialect; the Science DMZ is
established with the understanding that its specific gear may
differ from the rest of the estate and that this is acceptable.
The configuration is brought under version control across vendors.

\paragraph{What does not change.} The vendor mix itself is not
changed by this wave. Organizational-unit-level network autonomy,
where it exists, is respected; the wave provides the institutional
contract
under which the autonomy operates, not a new central authority.

\paragraph{Indicative duration.} Twelve to twenty-four months,
similar to the single-vendor trajectory but with a higher initial
investment in the cross-vendor configuration management layer and
a correspondingly lower marginal cost per subsequent zone.

%-----------------------------------------------------------------------------
\section{Reversibility criteria}
\label{sec:w3-reversibility}

Wave~3 is reversible at the per-segment level for an indicative
window of three to six months after each segment's cut-over.
Beyond that window, reversal is operationally possible but is
treated as a new project rather than as a wave rollback.

\begin{itemize}[leftmargin=1.5em,nosep]
  \item Each segment's pre-cut-over configuration is preserved in
    version control and is re-deployable if the post-cut-over
    state proves unworkable for that segment.
  \item Inter-segment routing rules established during coexistence
    are documented per segment, so that a reversal of one segment
    does not affect others.
  \item Where a vendor-specific standard interface proves
    operationally insufficient, the wave's response is documented
    --- substitution at next refresh, retention with documented
    sovereignty cost, or partial retention --- rather than a
    full-wave reversal.
\end{itemize}

%-----------------------------------------------------------------------------
\section{Exit criteria}
\label{sec:w3-exit}

Wave~3 closes when the four Network tests of the conformance suite
of \trchap{ch:conformance}{} (\S\ref{sec:cf-suite}) pass on a
representative sample of the segmentation, when the agreed segment
coverage is achieved, and when the Science DMZ (where present)
operates against its measurement targets.

\begin{itemize}[leftmargin=1.5em,nosep]
  \item NETCONF/YANG configuration apply succeeds without recourse
    to the supplier's management console; RADIUS attachment yields
    expected zone assignment; IPFIX flow records carry the
    canonical \texttt{correlation\_id}; configuration export and
    partial re-import on alternative-vendor switch restores
    expected connectivity on a test segment. (Network conformance
    subset.)
  \item Segment coverage: the institutionally agreed proportion of
    segments has been cut over. The remainder is in a documented
    schedule.
  \item Zone responsibility charters: every cut-over zone carries
    its charter, naming the accountable owner of the zone's
    published security profile.
  \item Science DMZ: where present, operating against its
    documented throughput and latency targets, with measurement
    data routed to the observability layer.
\end{itemize}

\begin{realworldbox}[Network segmentation precedents]
\noindent The Science DMZ pattern, developed at ESnet (Lawrence
Berkeley National Laboratory) and adopted at a number of
research-intensive universities, is one of the most extensively
documented patterns in academic network architecture. Multi-zone
campus segmentation is reported in various forms across large
research-intensive universities, among the institutions whose
architectural choices have been discussed in academic and
\gls{nren} community fora. At regional and
global scale the equivalent function is met by \gls{nren}- and
experiment-operated dedicated networks --- the regional and national
research and education networks' dedicated
circuits~\cite{geant} and the
CERN LHCOPN/LHCONE fabric~\cite{cern-lhcone} --- rather than
by a campus enclave. Specific deployment details should be verified
against current public documentation when cited.
\end{realworldbox}

%-----------------------------------------------------------------------------
\section{Regulatory anchoring}
\label{sec:w3-regulatory}

\noindent Wave~3 establishes zone-based segmentation and
attachment-time enforcement. It contributes to
\gls{gdpr}~Art.~32(1) (security through technical measures
applied at the network layer) and operationalises NIS2
Art.~21(2)(a) (risk-analysis applied to network design) and
(g) (cyber-hygiene baseline including segregation). It is the
principal anchor of 27001~A.8.20 (network security) and A.8.22
(segregation of networks) and contributes to NIST CSF
\textsf{PROTECT} (PR.IR --- infrastructure resilience). The
Science DMZ pattern operationalises segregation specifically for
high-throughput research traffic. Detailed mapping in
\trchap{ch:regulatory-mapping}{} \S\ref{sec:rm-gdpr},
\S\ref{sec:rm-nis2}, \S\ref{sec:rm-iso},
\S\ref{sec:rm-nist-csf}.
